\documentclass[10pt]{article}
\usepackage[margin=20mm]{geometry}
\usepackage{amsmath,amssymb,mathtools,fontspec,unicode-math,xurl,hyperref}
\IfFontExistsTF{DejaVu Serif}{\setmainfont{DejaVu Serif}}{\setmainfont{Latin Modern Roman}}
\setmathfont{latinmodern-math.otf}
\setlength{\parindent}{0pt}\setlength{\parskip}{6pt}\sloppy
\title{Technical Verification of
Understanding the training of infinitely deep and wide ResNets with Conditional Optimal Transport}
\author{Initial scientific review: M25-0418}
\date{29 October 2026}
\begin{document}
\maketitle
Abstract: The reviewed source is arXiv:2403.12887v2. Two local statements have incorrect orientations: the characteristic velocity in Proposition 3.2 and the coupling order in the first limit of Lemma 2.2. Explicit one-dimensional examples satisfy the stated hypotheses and contradict the printed formulas. The corresponding corrections are checked, and their implications are separated from the main convergence argument. Selected well-posedness hypotheses and the local Polyak–Łojasiewicz mechanism were also examined. The assessment is Minor for the demonstrated scope. Scientific adjudication and independent acceptance remain pending.\par
\section*{Identity and scope}
Raphaël Barboni, Gabriel Peyré and François-Xavier Vialard, Understanding the training of infinitely deep and wide ResNets with conditional optimal transport, Communications on Pure and Applied Mathematics 78 (2025), 2149–2205, DOI https://doi.org/10.1002/cpa.70004. The acquired source is the 49-page preprint arXiv:2403.12887v2, dated 21 July 2025 [BPV]. No identity with the publisher’s version of record was established. All page references below refer to the acquired PDF and agree with its printed numbering.\par
All source text was read, including the bibliography. This does not mean that every proof or cited external work was independently verified. The checks concentrate on the conditional transport conventions, the gradient-flow definition, the two statements below, the assumptions separating existence and uniqueness, and the kernel criterion for local convergence. No author-code execution or numerical training experiment is claimed.\par
Write \(s\in[0,1]\) for depth and \(t\) for optimization time. Measures \(\mu\in\mathcal P^{\mathrm{Leb}}_2([0,1]\times\Omega)\) have fixed first marginal \(ds\), and the conditional transport metric is \[d(\mu,\nu)^2=\int_0^1 W_2(\mu(\cdot\mid s),\nu(\cdot\mid s))^2\,ds.\] The ordered coupling convention on p. 9 assigns the first weight marginal to the first measure. We use the risk normalization in (5), \(L(\mu)=\mathbb E_{x,y}\ell(x_\mu(1),y)\), and denote the positive first variation in (17) by \(G_\mu=\nabla L[\mu]\). Definition 3.2 explicitly specifies velocity \(v_t=-G_{\mu_t}\).\par
\section*{Problem 1: the characteristic field has the wrong sign}
\subsection*{Claim and hypotheses}
Proposition 3.2 (p. 20) assumes that \(\phi\) is twice continuously differentiable and satisfies Assumptions 1–3. For a gradient flow \((\mu_t)_{t\geq0}\), it sets \[V_t(s,\omega)=(0,G_{\mu_t}(s,\omega)),\qquad
 \dot X_t=V_t\circ X_t,\quad X_0=\operatorname{Id},\] and asserts \(\mu_t=(X_t)_\#\mu_0\). The sign is confirmed by the actual source image; it conflicts with the continuity equation in Definition 3.2.\par
\subsection*{Check and result}
Take \(d=p=1\), \(\Omega=\mathbb R\), \(\phi(\omega,x)=\omega\), data distribution \(\delta_{(0,0)}\), and loss \(\ell(z,y)=\tfrac12(z-y)^2\). The feature map is smooth, has the required growth, and satisfies all three local Lipschitz and derivative assumptions. The data distribution is compactly supported. Set \[a(t)=e^{-t},\qquad \mu_t=ds\otimes\delta_{a(t)}.\] The forward and adjoint equations give, for every \(s\in[0,1]\), \[x_{\mu_t}(s)=s e^{-t},\qquad
 p_{\mu_t}(s)=e^{-t},\qquad G_{\mu_t}(s,\omega)=e^{-t}.\] This curve is absolutely continuous in \(d\), has finite second moment, and its actual velocity is \(a'(t)=-e^{-t}=-G_{\mu_t}\). It therefore is a gradient flow in the stated sense, for all \(t\geq0\).\par
The positive field prescribed by Proposition 3.2 instead sends the initial weight \(1\) to \[1+\int_0^t e^{-r}\,dr=2-e^{-t}.\] Thus its pushforward is \(ds\otimes\delta_{2-e^{-t}}\), and \[d\bigl((X_t)_\#\mu_0,\mu_t\bigr)=2(1-e^{-t})>0
 \quad(t>0).\] At \(t=\log2\) the distance is exactly \(1\). This is a counterexample to the printed representation, with no limiting argument or simulation.\par
\subsection*{Correction, counterevidence and impact}
The required field is \(V_t=(0,-G_{\mu_t})\). In the example its flow is \[X_t(s,\omega)=(s,\omega-1+e^{-t}),\] which has exactly the required pushforward. In general this sign matches the actual continuity equation, so the characteristic representation uses the negative field under the proposition’s regularity assumptions.\par
The source itself supplies counterevidence to a more serious interpretation. Definition 3.2 uses the minus sign, and Part 2 of Theorem 3.1 (p. 28) explicitly concludes \(v_t=-G_{\mu_t}\) from the energy-dissipation identity. The representation is cited for finite-width consistency and in Appendix A, but the latter’s support and uniqueness estimates use norms and remain compatible with the negative field. The example does not refute the local convergence theorem.\par
The finding is therefore Minor: an auxiliary identity is false as printed, while failure of a central conclusion has not been established. This classification is based on the demonstrated mathematical impact. The corrected example and sign are verified; no exhaustive reconstruction of the external characteristic-flow theorem or author adoption of this correction is claimed.\par
\section*{Problem 2: reversed coupling order in Lemma 2.2}
\subsection*{Claim and hypotheses}
Equation reference: eq:printed. The first assertion of Lemma 2.2 (p. 15) takes \(\gamma_t^h\in\Pi_o^{\mathrm{Leb}}(\mu_{t+h},\mu_t)\) and claims \[
 \left(\pi^1,\pi^2,\frac{\pi^3-\pi^2}{h}\right)_\#\gamma_t^h
 \longrightarrow (\pi^1,\pi^2,v_t)_\#\mu_t
 \quad\hbox{in }W_2 .\] Here \(v_t\) is the minimal tangent field from Proposition 2.4. The stated order of the measures and the displayed signs were checked on the source images.\par
\subsection*{Check and result}
Let \(I=(0,1)\), \(\Omega=\mathbb R\), and \(\mu_t=ds\otimes\delta_t\). Then \(d(\mu_t,\mu_u)=|t-u|\); the curve satisfies all the lemma’s assumptions and \(v_t=1\). For every interior \(t\) and sufficiently small nonzero \(h\), the unique declared conditional optimal coupling is \[\gamma_t^h=ds\otimes\delta_{(t+h,t)}.\] The left side of [eq:printed] (\#eq:printed) is consequently \(ds\otimes\delta_{(t+h,-1)}\), whereas the proposed limit is \(ds\otimes\delta_{(t,1)}\). Their ordinary \(W_2\) distance in the stated product space equals \[\sqrt{h^2+4}\longrightarrow 2.\] The claimed convergence fails at every such \(t\), so the almost-everywhere qualification does not remove the counterexample.\par
\subsection*{Correction, counterevidence and impact}
With the displayed pushforward unchanged, the coupling must be chosen in \(\Pi_o^{\mathrm{Leb}}(\mu_t,\mu_{t+h})\). For this example the corrected pushforward is exactly \(ds\otimes\delta_{(t,1)}\). Alternatively, keeping the printed coupling order requires the map \((\pi^1,\pi^3,(\pi^2-\pi^3)/h)\).\par
The proof on p. 16 says that the first marginal is \(\mu_t\) and uses the forward difference \(\varphi(\omega')-\varphi(\omega)\). Those steps agree with the corrected ordering. Moreover, the lemma’s second assertion, \[d\bigl(\mu_{t+h},(\operatorname{Id}+h(0,v_t))_\#\mu_t\bigr)=o(|h|),\] is exact in the translation example. This is the approximation used in the subsequent differentiation arguments. Thus the first display is false, but this example supplies no counterexample to the chain rule or the principal convergence result.\par
The finding is Minor. The correction is checked exactly in the example and matches the proof’s intended marginal convention. A complete independent proof of the cited general tangent-space theorem is outside this review. Journal correspondence and author adoption remain unverified.\par
\section*{Checks of the main framework}
\subsection*{Well-posedness and the descent convention}
Assumption 1 makes the forward NODE a Carathéodory equation with depth-integrable Lipschitz and growth bounds controlled by the second moment. Assumptions 2–3 support the local Lipschitz flow and adjoint estimates. With the stated risk convention, the chain rule and energy-dissipation identity give \[L'(\mu_t)=-\|G_{\mu_t}\|_{L^2(\mu_t)}^2.\] Conversely the maximal-slope argument on p. 28 gives \[-\langle v_t,G_{\mu_t}\rangle
 \geq \tfrac12\bigl(\|v_t\|^2+\|G_{\mu_t}\|^2\bigr),\] which forces \(v_t=-G_{\mu_t}\). This checks the core sign against Problem 1.\par
The hypotheses of Theorems 3.2–3.4 are distinct. The existence argument uses Assumption A with \(1\leq\alpha<2\); a uniform second-moment bound then controls the \(\alpha\)-moment tails needed in the narrow-limit argument. Uniqueness uses Assumption B on second derivatives and local Lipschitz continuity of the loss gradient. The finite-time stability estimate adds Assumption C and the quadratic loss. The general SHL architecture does not automatically satisfy all these global assumptions.\par
Appendix A treats compactly supported SHL initial measures. Its continuation argument can use a bound independent of an assumed global solution: along any local descent curve, \[d(\mu_t,\mu_0)\leq\sqrt{T L(\mu_0)},\qquad
 \mathcal E_2(\mu_t)\leq
 \bigl(\sqrt{\mathcal E_2(\mu_0)}+\sqrt{T L(\mu_0)}\bigr)^2
 \quad(0\leq t\leq T).\] These follow from energy dissipation and Cauchy–Schwarz. Together with the linear-growth characteristic estimate they supply a finite-time support bound for the cutoff construction. The apparent use of a full-orbit bound before existence therefore does not by itself establish an existence counterexample. The external minimizing-movement and transport theorems remain dependencies.\par
\subsection*{Local convergence and its scope}
For a curve of maximal slope satisfying \(|\nabla L|^2\geq mL\) in a ball of radius \(R\), its length before a first exit is at most \(2\sqrt{L(\mu_0)/m}\). The strict condition \(L(\mu_0)<mR^2/4\) thus prevents exit, gives \(L(\mu_t)\leq e^{-mt}L(\mu_0)\), and makes the tail length finite. This verifies the radius mechanism in Theorem 4.1, conditional on existence for the time interval in question.\par
For the normalization (5), the feature Gram matrix is positive semidefinite and the calculation on pp. 35–36 gives \[|\nabla L|^2
 =N^{-2}\int_0^1 p_\mu(s)^\top K_\mu(s)p_\mu(s)\,ds,\qquad
 \|p_\mu(s)\|^2\geq 2N e^{-C} L(\mu).\] The kernel lower bound therefore gives (31). The local Lipschitz estimate in Proposition 4.2 preserves half the initial eigenvalue integral \(\lambda_0\) on a ball of radius \(\min\{1,c(\mathcal E_2(\mu_0))\lambda_0/N\}\). Combining this with \(m\) proportional to \(\lambda_0/N\) explains the \(N^{-3}\lambda_0^3\) threshold in Theorem 4.2. The displayed risk on p. 35 has an extra factor \(1/2\) relative to (5), and the initial-loss value on p. 42 omits the half in its defined quadratic loss. The checks here use the explicit definition (5); these local normalization inconsistencies do not provide a further counterexample to the stated qualitative convergence criterion.\par
The theorem requires a positive initial Gram bound and sufficiently small initial risk, and concerns an existing gradient flow. The smooth SHL assumptions are not a well-posedness theorem for ReLU. Finite-feature conclusions use convergence in \(W_2\) and openness of the initialization conditions; no unconditional guarantee for arbitrary random initialization or discrete training is inferred. In Proposition 4.5, lifting makes the initial output in the target coordinates zero, while the rescaling increases the loss’s PL constant by \(\alpha^2\). This verifies the stated large-scaling mechanism when its Gram and smoothness assumptions hold.\par
\section*{Conclusion and evidence}
Dimension | Recorded state\par
Assessment | Minor: two supported local formula defects.\par
Repair verification | Corrected signs/orders checked; scope stated above.\par
Central conclusions | No failure established by these counterexamples.\par
Version correspondence | Publisher-final equality unverified.\par
Author adoption | Unknown.\par
Scientific adjudication | Pending; human reviews empty.\par
Independent acceptance | Not performed.\par
\section*{Supporting checks}
The exact source, full coverage ledger and current build records are retained under the numeric paper directory. The examples above require no author-supplied code. The complete external proofs, quantitative random-feature rates and all informal architecture comparisons were not independently verified.\par
No usable substantive predecessor report was found in the accessible preserved inventory. Existing source extraction and render bytes were reused with their historical attribution; current scientific checks and actual image views are recorded separately. No conclusion is drawn about inaccessible historical chat contents. The primary executor authored the report and inspected its current pages. Actual product-route information and runtime self-description are kept distinct in the provenance record.\par
1 R. Barboni, G. Peyré and F.-X. Vialard, Understanding the training of infinitely deep and wide ResNets with Conditional Optimal Transport, arXiv:2403.12887v2, 21 July 2025, https://arxiv.org/abs/2403.12887v2. Journal DOI: https://doi.org/10.1002/cpa.70004.\par
\end{document}
